\chapter{Hardware}

A hardware wallet is a small computer that holds keys
and shows you what you are about to sign. Used properly,
it is the last line when the laptop is dirty. Used
badly, it is an expensive way to rubber-stamp a phishing
site.

The standalone ``Hardware Wallets'' page on the site is
still a stub. The working guidance lives in the
intermediate-funds page. That is what this chapter
follows.

\section{Buy it like it matters}

Buy direct from the manufacturer. Not a marketplace.
Not a conference booth giveaway. Not a reseller with a
deal. Use a name that is not your public one if you can,
and a delivery point that is not your front door.

It must arrive uninitialized. No PIN. No seed card
already filled in. No ``we set it up for you.'' Reject
that device. Do not debug it. Do not ``just reset it.''

Then verify authenticity with the manufacturer's own
procedure before you generate anything. Ledger, Trezor,
GridPlus, and the others each have one. Use theirs.

\section{Set it up once, slowly}

Update firmware before you create keys. Set a PIN that
is not your phone PIN. Longer is better. Generate the
seed on the device. Do not import a seed from a
computer.

Do a dry-run recovery with the device's own feature
before you send funds. If you use a 25th-word
passphrase, check it in a temporary session and confirm
the address matches.

\section{What to look for}

A screen large enough to show the full destination and
the amount. Clear signing when the vendor has it. PIN
lockouts. A secure element with attestation if you are
buying that class of device. A vendor who publishes what
broke last time.

Clear signing is a help. It is not a substitute for
reading the screen. If the laptop and the device
disagree, the laptop is lying. Reject.

\section{A second device}

For anything you would hate to recover from a seed in a
hurry, keep a second hardware wallet with the same seed,
tested, stored separately. Confirm both still sign.
Monthly is enough if you actually do it.

Teams should not all buy the same brand. One firmware
bug should not take a quorum.
